\documentclass{overview}
\title{Calculus Overview - Derivative}

\begin{document}
\maketitle

\section{Existence of Derivative}

\subsection{Inferences}

\begin{theorem} \label{thm:diff_mult_cont}
If $f(x)$ is differentiable at $x =  x_0$ and $g(x)$ is continuous at $x = x_0$
but not differentiable at $x = x_0$, then the following functions
 \[
	h(x) = f(x)g(x)
\] 
are differentiable at $x = x_0$ if and only if $f(x_0) = 0$.
\end{theorem}
\begin{theorem}[Direct inference of Theorem \ref{thm:diff_mult_cont}]
	
	if $g(x)$ is continuous at $x = x_0$ then
	\[
		f(x) = |x-x_0|g(x)
	\] 
	is differentiable at $x = x_0$ if and only if $g(x_0) = 0$.

	Moreover, the function $f(x) = (x-x_0)^{n}|x-x_0|$ is differentiable to order
	$n$ at $x = x_0$ but not differentiable to order $n+1$ at $x = x_0$.
\end{theorem}

\begin{theorem}
	When $\alpha > 0$ and $\beta > 0$, the function
	\[
	f(x) = \begin{cases}
		x^{\alpha}\sin\dfrac{1}{x^{\beta}}, & x >0\\
		0, & x = 0 \\
		g(x), & x < 0
	\end{cases}
	\quad \text{or} \quad
	f(x) = \begin{cases}
		x^{\alpha}\cos\dfrac{1}{x^{\beta}}, & x >0\\
		0, & x = 0 \\
		g(x), & x < 0
	\end{cases}
	\] 
	where $g(x)$ satisfies $g(0^-) = 0, g'_-(0)=0$ and $g'(0^{-})=0$
	Then $f(x)$ at $x = 0$
	\begin{itemize}
		\item continuous if and only if $\alpha > 0$;
		\item differentiable if and only if $\alpha > 1$;
		\item has a derivative that is continuous if and only if
			$\alpha - \beta > 1$;
	\end{itemize}
	\begin{proof}
		The proof is direct from the definitions.
		\begin{itemize}
			\item continuity
				\[
					\lim_{x \to 0^{+}} f(x) = \lim_{x \to 0^{+}} x^{\alpha}
					\sin\frac{1}{x^{\beta}} = 0 = f(0) \iff \alpha > 0
				\] 
			\item differentiatility
				\[
					f'_+(0) = \lim_{x \to 0^{+}} \frac{f(x) - f(0)}{x - 0} =
					\lim_{x \to 0^{+}} x^{\alpha - 1} \sin\frac{1}{x^{\beta}}
					= 0 = f'_-(0) = g'_-(0) \iff \alpha > 1
				\] 
			\item continuity of derivative
				for $x > 0$,
				\[
					f'(x) = \alpha x^{\alpha - 1} \sin\frac{1}{x^{\beta}} -
					\beta x^{\alpha - \beta - 1} \cos\frac{1}{x^{\beta}}
				\] 
				therefore,
				\[
					\lim_{x \to 0^{+}} f'(x) = 0 = f'(0) \iff \alpha - \beta > 1
				\]
		\end{itemize}
		As for the case with cosine function, the proof is similar.
	\end{proof}
	\begin{remark}
		In fact, the $x^{\alpha}$ term can be replaced by $|x|^{\alpha}$ without
		affecting the conclusions. But we don't give the proof here.
	\end{remark}
\end{theorem}

\begin{theorem}
	When $f'_+(x_0)$ and $f'_-(x_0)$ both exist but are not equal, then
	$f(x)$ is not differentiable at $x = x_0$ but continuous at $x = x_0$.
	\begin{proof}
		\[
			f'_+(x_0) = \lim_{h \to 0^{+}} \frac{f(x_0 + h) - f(x_0)}{h} \text{ exists }
			\Rightarrow \lim_{h \to 0^{+}} f(x_0 + h) = f(x_0)
		\] 
		gives the right hand continuity at $x = x_0$, and similarly for the left hand side.
	\end{proof}
\end{theorem}

\begin{theorem}
	When $g(x)$ is n times differentiable at $x = x_0$ and $g(x_0) \neq 0$,
	then the function
	\[
		f(x) = (x-x_0)^{n}g(x)
	\] 
	has the following properties at $x = x_0$:
	\[
		f^{(k)}(x_0) = 0, \quad k = 0,1,2,\ldots,n-1
	\] 
\end{theorem}
\begin{theorem}
	\[
	f(x) = (x-x_1)^2 (x-x_2)^2 \cdots (x-x_n)^2
	\] 
	has $2n-1$ extreme points and  $2n-2$ inflection points.
\end{theorem}

\subsection{Common mistakes}

\begin{itemize}
	\item The left(or right) hand derivative not always equal to the left(or right) 
		hand limit of the derivative function. 

		\begin{example}
			Consider the function
			\[
				f(x) = \begin{cases}
					x^2\sin\frac{1}{x}, & x \neq 0 \\
					0, & x = 0
				\end{cases}
			\]
			We have
			\[
				f'(0) = \lim_{h \to 0} \frac{f(h) - f(0)}{h} = \lim_{h \to 0} h \sin\frac{1}{h} = 0
			\]
			but
			\[
				\lim_{x \to 0} f'(x) = \lim_{x \to 0} \left( 2x\sin\frac{1}{x} - \cos\frac{1}{x} \right)
			\]
			does not exist.
		\begin{remark}
		This mistake usually caused by assuming that the derivative function is
		continuous at that point, in which case the two limits are equal indeed.
		\end{remark}
		\end{example}
	\item Taking the limit before resolving the variable dependent part of a piecewise
		function.
		\begin{example}
			Try to determine whether the following function is differentiable at
			$x = 0$
			\[
				f(x) = \begin{cases}
					x, & x \leq 0 \\
					\frac{1}{n}, & \frac{1}{n+1} < x \leq \frac{1}{n}, n = 1,2,3,\ldots \\
				\end{cases}
			\]
			No doubt the left hand derivative at $x = 0$ is $1$. However,
			a common mistake is to think that the right hand derivative is $0$ since
			the graph is a horizontal line when approaching from the right side.
			But in fact, we have
			\[
				f'_{+}(0) = \lim_{x \to 0^{+}} \frac{f(x) - f(0)}{x} =
				\lim_{n \to \infty} \frac{\frac{1}{n}}{x}
			\] 
			where $1 / (n+1) < x \leq 1/n$ tells us that
			\[
				1= \frac{1 / n}{1 / n} \leq \frac{1 / n}{x} < \frac{1 / n}{1 / (n+1)}
				\to 1
			\] 
			as $n \to \infty$. Therefore, $f'_{+}(0) = 1$ and $f'(0) = 1$.
			
			The problem is that we cannot take the limit $x \to 0^{+}$ to get
			the function behavior near $0$ directly since the function is piecewise defined
			with variable dependent intervals.
			\begin{remark}
			This kind of mistake usually happens when the piecewise function has infinitely
			many pieces.
			\end{remark}
		\end{example}
	\item Has tangent line does not imply differentiability.
		\begin{example}
			Consider the function
			\[
				f(x) = \sqrt{x}
			\]
			at $x = 0$. The graph has a tangent line $x = 0$ at this point, but
			$f'(0)$ does not exist.
		\end{example}
	\item 0 is not necessarily a extreme point of even functions.
		Consider constant functions.
\end{itemize}

\section{Techniques of Derivative}
\subsection{Derivative of Products}

When only required to find the derivative at a certain point, it is not necessary
to actually find the derivative function.

One useful method is to factor the function such that one factor becomes zero at that point.
Then use the product rule to find the derivative at that point.
\begin{example}
	Consider the function
	\[
		f(x) = \frac{(x-1)(x-2)(x-3)\cdots(x-n)}{(x+1)(x+2)(x+3)\cdots(x+n)}
	\] 
	Find $f'(1)$.
	\begin{answer}
		Let $f(x) = (x-1)g(x)$, where
		\[
			g(x) = \frac{(x-2)(x-3)\cdots(x-n)}{(x+1)(x+2)(x+3)\cdots(x+n)}
		\]
		since
		\[
			f'(x) = (x-1)g'(x) + g(x)
		\] 
		we have $f'(1) = g(1)$. Therefore,
		\[
			f'(1) = \frac{(1-2)(1-3)\cdots(1-n)}{(1+1)(1+2)(1+3)\cdots(1+n)} =
			\frac{(-1)^{n-1}(n-1)!}{(n+1)!} = \frac{(-1)^{n-1}}{n(n+1)}
		\]
	\end{answer}
\end{example}

Another useful method is to use the definition of derivative directly to
find the derivative at that point.
\begin{example}
	Consider the function
	\[
		f(x) = (e^{x}-1)(e^{2x}-2)(e^{3x}-3)\cdots(e^{nx}-n)
	\]
	Find $f'(0)$.
	\begin{answer}
		Using the definition of derivative, we have
		\[
		\begin{aligned}
			f'(0) &= \lim_{h \to 0} \frac{f(h) - f(0)}{h} =
			\lim_{h \to 0} \frac{(e^{h}-1)(e^{2h}-2)(e^{3h}-3)\cdots(e^{nh}-n)}{h} \\
						&= \lim_{h \to 0} \frac{e^{h}-1}{h}\cdot (e^{2h}-2)(e^{3h}-3)\cdots(e^{nh}-n) \\
						&= 1 \cdot (-2)(-3)\cdots(-n) = (-1)^{n-1} n!
		\end{aligned}
	\]
	\end{answer}
\end{example}

\begin{example}
	Consider the function
	\[
		f(x) = \frac{\ln|x|}{2x^2-\ln|x|}
	\]
	Find $f'(-1)$.
	\begin{answer}
		Using the definition of derivative, we have
		\[
		\begin{aligned}
			f'(-1) &= \lim_{x \to -1} \frac{\ln|x|}{(x+1)(2x^2 - \ln|x|)} \\
					&= \lim_{x \to -1} \frac{-x-1}{(x+1)(2x^2 - \ln|x|)} \\
					&= -\frac{1}{2}
		\end{aligned}
		\]
	\end{answer}
\end{example}

Another method is to use logarithmic differentiation.
\begin{example}
	Consider the function
	\[
		\varphi(x) = \frac{\exp(x)\exp(f_1(x)f_2(x))}{f_1(x)f_2(x)}
	\]
	where $f_1(x)$ and $f_2(x)$ are differentiable functions and positive in 
	their domains. Also satisfy $f_1'(0)=f_2'(0)=0$.
	Try to prove that $\varphi'(0) > 1$
	\begin{answer}
		\[
		\begin{aligned}
			\ln \varphi(x) &= x + f_1(x)f_2(x) - \ln f_1(x) - \ln f_2(x) \\
			\Rightarrow \quad \varphi'(x) &= \varphi(x) \left[ 1 +
			f_1'(x)f_2(x) + f_1(x)f_2'(x) - \frac{f_1'(x)}{f_1(x)} -
			\frac{f_2'(x)}{f_2(x)} \right]
		\end{aligned}
		\] 
		therefore
		\[
			\varphi'(0) = \varphi(0) = \frac{\exp(f_1(0)f_2(0))}{f_1(0)f_2(0)} > 1
		\]
	\end{answer}
\end{example}

\subsection{Derivative of Implicit Functions}

\begin{theorem}
	Given a function $x=f^{-1}(y)$ which is the inverse of $y=f(x)$. If $f(x)$ is
	2th order differentiable and $f'(x) \neq 0$, then the second order derivative of
	the inverse function is given by
	\[
		\frac{\dd^2 x}{\dd y^2} = -\frac{f''(x)}{[f'(x)]^3}
	\] 
	\begin{proof}
		\[
			\frac{\dd^2 x}{\dd y^2} = \frac{\dd \frac{\dd x}{\dd y} / \dd x}{\dd y / \dd x} =
			\frac{\dd \frac{1}{f'(x)}}{\dd x} \cdot \frac{1}{f'(x)} =
			-\frac{f''(x)}{[f'(x)]^3}
		\] 
	\end{proof}
\end{theorem}
\begin{theorem}
	Given a function $y=f(x)$ which is defined implicitly by the parametric equations
	$\begin{cases} x = x(t)\\ y = y(t)\end{cases}\!\!\!\!\!\!$,then 
	\[
		\frac{\dd^2 y}{\dd x^2} = \frac{x'(t)y''(t) - y'(t)x''(t)}{[x'(t)]^3}
	\] 
	\begin{remark}
		We don't need to remember this formula, just remember how to derive it.
	\end{remark}
\end{theorem}

\subsection{Trick to remember Derivative of $f(x)^{g(x)}$ }
We have an formula but hard to remember:
\[
	\frac{\dd}{\dd x} \left[ f(x)^{g(x)} \right] =
	f(x)^{g(x)} \left[ g'(x) \ln f(x) + g(x) \frac{f'(x)}{f(x)} \right]
\] 
A simple trick is to rewrite it as
\[
	\frac{\dd}{\dd x} \left[ f(x)^{g(x)} \right] =
	f(x)^{g(x)} \frac{\dd}{\dd x} \left[ g(x) \ln f(x) \right]
\] 
or more simply,
\[
	(u^{v})' = u^{v} (\ln(u^{v}))'
\] 
\begin{example}
	Find the derivative of $y = x^{x^x}$.
	\begin{answer}
		we have
		\[
		\begin{aligned}
			y' &= x^{x^x} \frac{\dd}{\dd x} \left[ x^x \ln x \right] \\
				 &= x^{x^x} \left[ x^x (x\ln x)'\ln x + x^{x-1} \right] \\
				 &= x^{x^x} \left[ x^x (\ln x + 1) \ln x + x^{x-1} \right]
		\end{aligned}
		\]
	\end{answer}
\end{example}

\section{Higher Order Derivatives}

here are some higher order derivative of common functions:

\begin{itemize}
	\item $(e^{ax+b})^{(n)} = a^{n} e^{ax+b}$
	\item $(a ^{x})^{(n)} = (\ln a)^{n} a^{x}$
	\item If $m$ is a positive integer,
		 \[
			 (x^m)^{(n)} = \begin{cases}
				 m(m-1)(m-2)\cdots(m-n+1)x^{m-n}, & m > n\\
				 n! , & m = n \\
				 0, & m < n
			 \end{cases}
		\] 
	\item If $\alpha$ is any real number, in general,
		 \[
			 [(ax + b)^{\alpha}]^{(n)} = \alpha(\alpha - 1)(\alpha - 2)\cdots
			 (\alpha - n + 1)(ax + b)^{\alpha - n}a^{n}
		\]
		Especially, when $\alpha = -1$,
		 \[
			 \left( \frac{1}{x} \right)^{(n)} =
			 \frac{(-1)^{n} n!}{x^{n+1}}
		\] 
	\item For logarithmic functions, is similar to the power functions
		\[
		[\ln(x)]^{(n)} = \dfrac{(-1)^{n-1}(n-1)!}{x^{n}}
		\] 
	\item For trigonometric functions,
		\[
		(\sin x)^{(n)} = \sin\left( x + \frac{n\pi}{2} \right), \quad
		(\cos x)^{(n)} = \cos\left( x + \frac{n\pi}{2} \right)
		\]
\end{itemize}

\begin{example}
	Find the $n$th derivative of $y = \dfrac{1-x}{1+x}$.
	\begin{answer}
		Rewrite $y$ as
		\[
			y = \frac{2}{1+x} - 1
		\]
		therefore,
		\[
			y^{(n)} = 2 \cdot \frac{(-1)^{n} n!}{(1+x)^{n+1}}
		\]
	\end{answer}
\end{example}

The Leibniz formula for the $n$th derivative of a product of two functions
\[
	(uv)^{(n)} = \sum_{k=0}^{n} C_n^k u^{(k)} v^{(n-k)}
\] 
Especially, we have
\[
	(xe^{x})^{(n)} = (x+n)e^{x}
\] 
And when $u = v$, we have
\[
	(u^2)^{(n)} = \sum_{k=0}^{n} C_n^k u^{(k)} u^{(n-k)}
\]

Taylor series expansion can also be used to find higher order derivatives
at a certain point.
\[
	f^{(n)}(x_0) = n!a_n
\] 
where $a_n$ is the coefficient of $(x - x_0)^{n}$ in the Taylor series expansion
of $f(x)$ at $x = x_0$.

\subsection{Recursive method to find Higher Order Derivatives}
First, calculate the first few order derivatives, and try to find a relation
between the lower order and original function. Then evaluate nth order derivative
using the relation.
\begin{example}
	Find the $n$th derivative at $x=0$ of $y = \arcsin x$.
	\begin{answer}
		First, we have
		\[
			(\arcsin x)' = \frac{1}{\sqrt{1-x^2}}
		\]
		gives
		\[
			(1-x^2)(y')^2 = 1
		\]
		to get rid of the square term. Then differentiate both sides,
		\[
		\begin{aligned}
			-2x(y')^2 + 2(1-x^2)y'y'' &= 0
			\Rightarrow (1-x^2)y''- xy' = 0
		\end{aligned}
		\]
		take nth derivative on both sides again,
		\[
			(1-x^2)y^{(n+2)} - 2n x y^{(n+1)} - n(n-1) y^{(n)} - x y^{(n+1)} - n y^{(n)} = 0
		\]
		evaluate at $x = 0$ gives
		\[
			y^{(n+2)}(0) = n^2 y^{(n)}(0)
		\]
		Using the initial values $y^{(1)}(0) = 1, y^{(2)}(0) = 0$, we have
		\[
			y^{(n)}(0) = \begin{cases}
				[(n-2)!!]^2, & n \text{ is odd} \\
				0, & n \text{ is even}
			\end{cases}
		\]
	\end{answer}
\end{example}

\section{Rolle's Theorem on $f'+gf=0$}

Here we focus on a class of problems can be converted to proof the existence
of roots of the equation $f'(x) + g(x) f(x) = 0$. More specifically, like proving
that there exists at least one point $c$ such that $f'(c)=-f(c) / c$ or $f'(c)+\lambda f(c)=0$

To prove the existence of such point, we want to construct a function $F(x)$ such that
\[
	F'(x) = 0 \iff f'(x) + g(x) f(x) = 0 \quad \text{and} \quad F(a) = F(b)
\] 
Then we can use Rolle's theorem to prove the existence of such point.

A simple idea is to let $F(x)$ be the integration of $f'(x) + g(x) f(x)$ directly.
However, it is not easy to find the result. Instead, we use a method called
``integrating factor'' in differential equations to construct $F(x)$.

we want to find a function $\mu(x)$ such that when we let
\[
	F(x) = \mu(x) f(x)
\]
then $F'(x)$ can be expressed as
\[
	F'(x) = \mu(x) [f'(x) + g(x) f(x)]
\]
To do this, we need to have
\[
	F'(x) = \mu'(x) f(x) + \mu(x) f'(x) = \mu(x) [f'(x) + g(x) f(x)]
\]
which gives
\[
	\mu'(x) = g(x) \mu(x)
\]
that is
\[
	\mu(x) = e^{\int g(x) \dd x}
\]
then test whether $F(a) = F(b)$ holds, if so, then by Rolle's theorem we can
conclude what we want.

In fact, this method also works with below form:
\begin{theorem}
	\[
		e^{\int g(x) \dd x} [f'(x)+g(x)f(x)+p(x)] = \left[f(x)
		e^{\int g(x) \dd x} + \int p(x) e^{\int g(x) \dd x} \dd x \right]'
	\] 
\end{theorem}
\begin{example}
	Prove that there exists at least one point $c \in (a,b)$ such that
	\[
		f'(c) + \frac{f(c)}{c} = 0
	\] 
	where $f(x)$ is differentiable on $[a,b]$ and $f(a) = f(b) = 0$,
	$0 < a < b$.
	\begin{answer}
		Let
		\[
			F(x) = f(x) e^{\int \frac{1}{x} \dd x} = x f(x)
		\] 
		then
		\[
			F'(x) = e^{\int \frac{1}{x} \dd x} [f'(x) + \frac{f(x)}{x}] =
			x [f'(x) + \frac{f(x)}{x}]
		\] 
		since $F(a) = a f(a) = 0$ and $F(b) = b f(b) = 0$, by Rolle's theorem,
		there exists at least one point $c \in (a,b)$ such that $F'(c) = 0$,
		that is
		\[
			f'(c) + \frac{f(c)}{c} = 0
		\]
	\end{answer}
\end{example}

\section{Asymptotes}

When $\lim_{x \to +\infty}\frac{ f(x)}{x}$ exists it does not necessarily mean that
there is an oblique asymptote. For example, consider the function
\[
	f(x) = x + \sin x
\]
we have
\[
	\lim_{x \to +\infty} \frac{f(x)}{x} = 1
\]
but there is no oblique asymptote since 
\[
	\lim_{x \to +\infty} [f(x) - x] = \lim_{x \to +\infty} \sin x
\]
does not exist.

\section{Curvature}

The curvature of a curve at a certain point is defined as
\[
	\kappa = \left| \frac{\dd \alpha}{\dd s} \right| = 
	\frac{|y''|}{[1 + (y')^2]^{3/2}}
\]
and the radius of curvature is defined as
\[
	\rho = \frac{1}{\kappa} = \frac{[1 + (y')^2]^{3/2}}{|y''|}
\]


\end{document}
